\section{详细设计}
\subsection{本章概述}
详细设计重点说明系统如何通过 SQL 支撑核心业务查询和统计分析。以下查询围绕访客预约、审批流转、门岗通行、风险控制和统计报表展开，体现关系数据库在数据关联、聚合统计和状态筛选方面的设计作用。

\subsection{SQL-01 查询某访客的历史预约记录}
\textbf{查询名称：}查询某访客的历史预约记录。\par
\textbf{查询目的：}用于访客记录查询页，按手机号回溯访客所有历史预约，关联展示被访部门、被访人、预约状态和访问状态。\par
\textbf{涉及数据表：}\texttt{visit\_apply, visitor, department, sys\_user}。\par
\textbf{SQL 语句：}
\begin{lstlisting}[style=sqlstyle]
SELECT
  v.visitor_name,
  v.phone,
  va.apply_no,
  va.visit_reason,
  d.dept_name,
  u.real_name AS host_name,
  va.plan_start_time,
  va.plan_end_time,
  va.apply_status,
  va.access_status
FROM visit_apply va
JOIN visitor v ON va.visitor_id = v.id
LEFT JOIN department d ON va.department_id = d.id
LEFT JOIN sys_user u ON va.host_user_id = u.id
WHERE va.deleted = 0 AND v.phone = '13950000001'
ORDER BY va.plan_start_time DESC;
\end{lstlisting}
\textbf{查询说明：}该查询与系统页面、后端统计接口或课程设计报告中的业务说明保持一致，能够直接基于当前数据库脚本执行。

\subsection{SQL-02 查询某被访人的待确认预约}
\textbf{查询名称：}查询某被访人的待确认预约。\par
\textbf{查询目的：}用于被访人端待确认列表，保证被访人只处理与本人账号相关的待确认预约。\par
\textbf{涉及数据表：}\texttt{visit\_apply, visitor, sys\_user}。\par
\textbf{SQL 语句：}
\begin{lstlisting}[style=sqlstyle]
SELECT va.apply_no, v.visitor_name, v.phone, va.visit_reason, va.plan_start_time, va.apply_status
FROM visit_apply va
JOIN visitor v ON va.visitor_id = v.id
JOIN sys_user host ON va.host_user_id = host.id
WHERE va.deleted = 0
  AND host.username = 'teacher01'
  AND va.apply_status = 'PENDING_HOST'
ORDER BY va.plan_start_time ASC;
\end{lstlisting}
\textbf{查询说明：}该查询与系统页面、后端统计接口或课程设计报告中的业务说明保持一致，能够直接基于当前数据库脚本执行。

\subsection{SQL-03 查询某部门的待审批预约}
\textbf{查询名称：}查询某部门的待审批预约。\par
\textbf{查询目的：}用于部门审批页，筛选本部门已由被访人确认或进入部门审批阶段的预约申请。\par
\textbf{涉及数据表：}\texttt{visit\_apply, visitor, department, sys\_user}。\par
\textbf{SQL 语句：}
\begin{lstlisting}[style=sqlstyle]
SELECT va.apply_no, d.dept_name, v.visitor_name, host.real_name AS host_name, va.visit_reason, va.plan_start_time, va.apply_status
FROM visit_apply va
JOIN visitor v ON va.visitor_id = v.id
JOIN department d ON va.department_id = d.id
LEFT JOIN sys_user host ON va.host_user_id = host.id
WHERE va.deleted = 0
  AND d.dept_name = '信息与通信工程学院'
  AND va.apply_status IN ('HOST_CONFIRMED', 'PENDING_DEPT')
ORDER BY va.plan_start_time ASC;
\end{lstlisting}
\textbf{查询说明：}该查询与系统页面、后端统计接口或课程设计报告中的业务说明保持一致，能够直接基于当前数据库脚本执行。

\subsection{SQL-04 查询当天已入校访客}
\textbf{查询名称：}查询当天已入校访客。\par
\textbf{查询目的：}用于门岗和统计页面，动态查询当天已经完成入校登记的访客。\par
\textbf{涉及数据表：}\texttt{access\_record, visit\_apply, visitor, campus\_gate, sys\_user}。\par
\textbf{SQL 语句：}
\begin{lstlisting}[style=sqlstyle]
SELECT ar.id, va.apply_no, v.visitor_name, v.phone, g.gate_name AS entry_gate, guard.real_name AS guard_name, ar.entry_time
FROM access_record ar
JOIN visit_apply va ON ar.apply_id = va.id
JOIN visitor v ON ar.visitor_id = v.id
LEFT JOIN campus_gate g ON ar.entry_gate_id = g.id
LEFT JOIN sys_user guard ON ar.entry_guard_id = guard.id
WHERE ar.deleted = 0
  AND ar.entry_time >= CURRENT_DATE()
  AND ar.entry_time < DATE_ADD(CURRENT_DATE(), INTERVAL 1 DAY)
ORDER BY ar.entry_time DESC;
\end{lstlisting}
\textbf{查询说明：}该查询与系统页面、后端统计接口或课程设计报告中的业务说明保持一致，能够直接基于当前数据库脚本执行。

\subsection{SQL-05 查询当天已入校但未离校访客}
\textbf{查询名称：}查询当天已入校但未离校访客。\par
\textbf{查询目的：}用于当前在校访客页，查询当天已入校但尚未离校的访客记录。\par
\textbf{涉及数据表：}\texttt{access\_record, visit\_apply, visitor, department}。\par
\textbf{SQL 语句：}
\begin{lstlisting}[style=sqlstyle]
SELECT ar.id, va.apply_no, v.visitor_name, v.phone, d.dept_name, ar.entry_time, ar.access_status
FROM access_record ar
JOIN visit_apply va ON ar.apply_id = va.id
JOIN visitor v ON ar.visitor_id = v.id
LEFT JOIN department d ON va.department_id = d.id
WHERE ar.deleted = 0
  AND ar.access_status = 'ENTERED'
  AND ar.entry_time >= CURRENT_DATE()
  AND ar.entry_time < DATE_ADD(CURRENT_DATE(), INTERVAL 1 DAY)
ORDER BY ar.entry_time DESC;
\end{lstlisting}
\textbf{查询说明：}该查询与系统页面、后端统计接口或课程设计报告中的业务说明保持一致，能够直接基于当前数据库脚本执行。

\subsection{SQL-06 查询超时未离校访客}
\textbf{查询名称：}查询超时未离校访客。\par
\textbf{查询目的：}用于超时未离校页面，结合访问状态与计划离校时间识别风险访客。\par
\textbf{涉及数据表：}\texttt{access\_record, visit\_apply, visitor}。\par
\textbf{SQL 语句：}
\begin{lstlisting}[style=sqlstyle]
SELECT va.apply_no, v.visitor_name, v.phone, va.plan_end_time, ar.entry_time, ar.overtime_flag, ar.abnormal_reason
FROM access_record ar
JOIN visit_apply va ON ar.apply_id = va.id
JOIN visitor v ON ar.visitor_id = v.id
WHERE ar.deleted = 0
  AND (ar.access_status = 'OVERTIME' OR (ar.exit_time IS NULL AND va.plan_end_time < NOW()))
ORDER BY va.plan_end_time ASC;
\end{lstlisting}
\textbf{查询说明：}该查询与系统页面、后端统计接口或课程设计报告中的业务说明保持一致，能够直接基于当前数据库脚本执行。

\subsection{SQL-07 按部门统计本月访客数量}
\textbf{查询名称：}按部门统计本月访客数量。\par
\textbf{查询目的：}用于部门访客排行图表，按本月访问申请统计各部门接待量。\par
\textbf{涉及数据表：}\texttt{visit\_apply, department}。\par
\textbf{SQL 语句：}
\begin{lstlisting}[style=sqlstyle]
SELECT d.dept_name, COUNT(*) AS visitor_count
FROM visit_apply va
JOIN department d ON va.department_id = d.id
WHERE va.deleted = 0
  AND va.plan_start_time >= DATE_FORMAT(CURRENT_DATE(), '%Y-%m-01')
  AND va.plan_start_time < DATE_ADD(DATE_FORMAT(CURRENT_DATE(), '%Y-%m-01'), INTERVAL 1 MONTH)
GROUP BY d.id, d.dept_name
ORDER BY visitor_count DESC;
\end{lstlisting}
\textbf{查询说明：}该查询与系统页面、后端统计接口或课程设计报告中的业务说明保持一致，能够直接基于当前数据库脚本执行。

\subsection{SQL-08 按校门统计本月通行数量}
\textbf{查询名称：}按校门统计本月通行数量。\par
\textbf{查询目的：}用于校门通行统计图表，分别统计本月各校门入校、离校和总通行次数。\par
\textbf{涉及数据表：}\texttt{campus\_gate, access\_record}。\par
\textbf{SQL 语句：}
\begin{lstlisting}[style=sqlstyle]
SELECT g.gate_name,
       SUM(CASE WHEN ar.entry_gate_id = g.id THEN 1 ELSE 0 END) AS entry_count,
       SUM(CASE WHEN ar.exit_gate_id = g.id THEN 1 ELSE 0 END) AS exit_count,
       SUM(CASE WHEN ar.entry_gate_id = g.id THEN 1 ELSE 0 END) + SUM(CASE WHEN ar.exit_gate_id = g.id THEN 1 ELSE 0 END) AS total_count
FROM campus_gate g
LEFT JOIN access_record ar ON ar.deleted = 0
  AND (ar.entry_gate_id = g.id OR ar.exit_gate_id = g.id)
  AND COALESCE(ar.entry_time, ar.exit_time) >= DATE_FORMAT(CURRENT_DATE(), '%Y-%m-01')
  AND COALESCE(ar.entry_time, ar.exit_time) < DATE_ADD(DATE_FORMAT(CURRENT_DATE(), '%Y-%m-01'), INTERVAL 1 MONTH)
WHERE g.deleted = 0
GROUP BY g.id, g.gate_name
ORDER BY total_count DESC;
\end{lstlisting}
\textbf{查询说明：}该查询与系统页面、后端统计接口或课程设计报告中的业务说明保持一致，能够直接基于当前数据库脚本执行。

\subsection{SQL-09 查询黑名单访客}
\textbf{查询名称：}查询黑名单访客。\par
\textbf{查询目的：}用于黑名单管理和预约风险检查，展示仍处于生效状态的限制入校人员。\par
\textbf{涉及数据表：}\texttt{blacklist, visitor}。\par
\textbf{SQL 语句：}
\begin{lstlisting}[style=sqlstyle]
SELECT b.id, v.visitor_name, b.phone, b.id_number, b.reason, b.level, b.status, b.start_time, b.end_time
FROM blacklist b
LEFT JOIN visitor v ON b.visitor_id = v.id
WHERE b.deleted = 0 AND b.status = 'ACTIVE'
ORDER BY b.level DESC, b.start_time DESC;
\end{lstlisting}
\textbf{查询说明：}该查询与系统页面、后端统计接口或课程设计报告中的业务说明保持一致，能够直接基于当前数据库脚本执行。

\subsection{SQL-10 查询某预约的完整审批记录}
\textbf{查询名称：}查询某预约的完整审批记录。\par
\textbf{查询目的：}用于预约详情页，按审批顺序展示某个预约的完整审批链路。\par
\textbf{涉及数据表：}\texttt{approval\_record, visit\_apply, sys\_user}。\par
\textbf{SQL 语句：}
\begin{lstlisting}[style=sqlstyle]
SELECT va.apply_no, ar.approval_step, approver.real_name AS approver_name, ar.approval_result, ar.approval_comment, ar.approval_time
FROM approval_record ar
JOIN visit_apply va ON ar.apply_id = va.id
LEFT JOIN sys_user approver ON ar.approver_user_id = approver.id
WHERE ar.deleted = 0 AND va.apply_no = (SELECT apply_no FROM visit_apply ORDER BY id LIMIT 1)
ORDER BY ar.sort_order ASC, ar.approval_time ASC;
\end{lstlisting}
\textbf{查询说明：}该查询与系统页面、后端统计接口或课程设计报告中的业务说明保持一致，能够直接基于当前数据库脚本执行。

\subsection{SQL-11 查询某访客的车辆和随行人员}
\textbf{查询名称：}查询某访客的车辆和随行人员。\par
\textbf{查询目的：}用于预约详情页，联合展示车辆和随行人员信息。\par
\textbf{涉及数据表：}\texttt{visit\_apply, visitor, visitor\_vehicle, visitor\_companion}。\par
\textbf{SQL 语句：}
\begin{lstlisting}[style=sqlstyle]
SELECT va.apply_no, v.visitor_name, vv.plate_no, vv.vehicle_type, vv.color, vc.companion_name, vc.phone AS companion_phone, vc.relation_remark
FROM visit_apply va
JOIN visitor v ON va.visitor_id = v.id
LEFT JOIN visitor_vehicle vv ON va.vehicle_id = vv.id
LEFT JOIN visitor_companion vc ON vc.apply_id = va.id AND vc.deleted = 0
WHERE va.deleted = 0 AND va.id = 2;
\end{lstlisting}
\textbf{查询说明：}该查询与系统页面、后端统计接口或课程设计报告中的业务说明保持一致，能够直接基于当前数据库脚本执行。

\subsection{SQL-13 查询今日访客概览}
\textbf{查询名称：}查询今日访客概览。\par
\textbf{查询目的：}用于首页驾驶舱，汇总今日访客、当前在校、超时未离校、待审批和黑名单风险。\par
\textbf{涉及数据表：}\texttt{visit\_apply, access\_record, blacklist}。\par
\textbf{SQL 语句：}
\begin{lstlisting}[style=sqlstyle]
SELECT
  (SELECT COUNT(*) FROM visit_apply WHERE deleted = 0 AND plan_start_time >= CURRENT_DATE() AND plan_start_time < DATE_ADD(CURRENT_DATE(), INTERVAL 1 DAY)) AS today_visitors,
  (SELECT COUNT(*) FROM access_record WHERE deleted = 0 AND access_status = 'ENTERED') AS current_visitors,
  (SELECT COUNT(*) FROM access_record WHERE deleted = 0 AND access_status = 'OVERTIME') AS overtime_visitors,
  (SELECT COUNT(*) FROM visit_apply WHERE deleted = 0 AND apply_status IN ('PENDING_HOST','HOST_CONFIRMED','PENDING_DEPT')) AS pending_approvals,
  (SELECT COUNT(*) FROM blacklist WHERE deleted = 0 AND status = 'ACTIVE') AS blacklist_risks;
\end{lstlisting}
\textbf{查询说明：}该查询与系统页面、后端统计接口或课程设计报告中的业务说明保持一致，能够直接基于当前数据库脚本执行。

\subsection{SQL-14 查询近七天访客趋势}
\textbf{查询名称：}查询近七天访客趋势。\par
\textbf{查询目的：}用于近七天访客趋势图，使用递归日期表保证没有访问的日期也能显示为 0。\par
\textbf{涉及数据表：}\texttt{days, visit\_apply}。\par
\textbf{SQL 语句：}
\begin{lstlisting}[style=sqlstyle]
WITH RECURSIVE days AS (
  SELECT DATE_SUB(CURRENT_DATE(), INTERVAL 6 DAY) AS stat_date
  UNION ALL
  SELECT DATE_ADD(stat_date, INTERVAL 1 DAY) FROM days WHERE stat_date < CURRENT_DATE()
)
SELECT DATE_FORMAT(days.stat_date, '%m-%d') AS stat_date,
       COUNT(va.id) AS visitor_count
FROM days
LEFT JOIN visit_apply va ON va.deleted = 0
  AND va.plan_start_time >= days.stat_date
  AND va.plan_start_time < DATE_ADD(days.stat_date, INTERVAL 1 DAY)
GROUP BY days.stat_date
ORDER BY days.stat_date;
\end{lstlisting}
\textbf{查询说明：}该查询与系统页面、后端统计接口或课程设计报告中的业务说明保持一致，能够直接基于当前数据库脚本执行。

\subsection{SQL-15 查询审批通过率}
\textbf{查询名称：}查询审批通过率。\par
\textbf{查询目的：}用于审批通过率统计，计算部门审批通过、拒绝/退回、待处理数量和通过率。\par
\textbf{涉及数据表：}\texttt{approval\_record}。\par
\textbf{SQL 语句：}
\begin{lstlisting}[style=sqlstyle]
SELECT
  SUM(CASE WHEN approval_result = 'APPROVED' THEN 1 ELSE 0 END) AS approved_count,
  SUM(CASE WHEN approval_result IN ('REJECTED','RETURNED') THEN 1 ELSE 0 END) AS rejected_count,
  SUM(CASE WHEN approval_result = 'PENDING' THEN 1 ELSE 0 END) AS pending_count,
  ROUND(SUM(CASE WHEN approval_result = 'APPROVED' THEN 1 ELSE 0 END) * 100.0 / COUNT(*), 2) AS pass_rate
FROM approval_record
WHERE deleted = 0 AND approval_step = 'DEPT_APPROVAL';
\end{lstlisting}
\textbf{查询说明：}该查询与系统页面、后端统计接口或课程设计报告中的业务说明保持一致，能够直接基于当前数据库脚本执行。

\subsection{SQL-18 访问事由分布}
\textbf{查询名称：}访问事由分布。\par
\textbf{查询目的：}用于访问事由分布图，通过关键词归类展示家长、企业合作、后勤维修、学术交流等访问目的。\par
\textbf{涉及数据表：}\texttt{visit\_apply}。\par
\textbf{SQL 语句：}
\begin{lstlisting}[style=sqlstyle]
SELECT reason_type, COUNT(*) AS reason_count
FROM (
  SELECT CASE
    WHEN visit_reason LIKE '%家长%' OR visit_reason LIKE '%学籍%' THEN '学生家长'
    WHEN visit_reason LIKE '%企业%' OR visit_reason LIKE '%合作%' THEN '企业合作'
    WHEN visit_reason LIKE '%维修%' OR visit_reason LIKE '%后勤%' THEN '后勤维修'
    WHEN visit_reason LIKE '%学术%' OR visit_reason LIKE '%研讨%' THEN '学术交流'
    WHEN visit_reason LIKE '%面试%' OR visit_reason LIKE '%就业%' THEN '面试招聘'
    WHEN visit_reason LIKE '%快递%' OR visit_reason LIKE '%配送%' THEN '快递配送'
    WHEN visit_reason LIKE '%校友%' THEN '校友返校'
    WHEN visit_reason LIKE '%专家%' OR visit_reason LIKE '%评审%' OR visit_reason LIKE '%讲座%' THEN '外聘专家'
    ELSE '其他'
  END AS reason_type
  FROM visit_apply
  WHERE deleted = 0
    AND plan_start_time >= DATE_FORMAT(CURRENT_DATE(), '%Y-%m-01')
    AND plan_start_time < DATE_ADD(DATE_FORMAT(CURRENT_DATE(), '%Y-%m-01'), INTERVAL 1 MONTH)
) t
GROUP BY reason_type
ORDER BY reason_count DESC;
\end{lstlisting}
\textbf{查询说明：}该查询与系统页面、后端统计接口或课程设计报告中的业务说明保持一致，能够直接基于当前数据库脚本执行。

\subsection{SQL-19 审批状态分布}
\textbf{查询名称：}审批状态分布。\par
\textbf{查询目的：}用于审批状态分布图，展示预约在待确认、审批通过、审批拒绝、取消等状态上的数量。\par
\textbf{涉及数据表：}\texttt{visit\_apply}。\par
\textbf{SQL 语句：}
\begin{lstlisting}[style=sqlstyle]
SELECT apply_status, COUNT(*) AS status_count
FROM visit_apply
WHERE deleted = 0
GROUP BY apply_status
ORDER BY status_count DESC;
\end{lstlisting}
\textbf{查询说明：}该查询与系统页面、后端统计接口或课程设计报告中的业务说明保持一致，能够直接基于当前数据库脚本执行。

\subsection{SQL-20 入校离校状态分布}
\textbf{查询名称：}入校离校状态分布。\par
\textbf{查询目的：}用于入校离校状态分布图，统计未入校、已入校、已离校、超时未离校等访问状态。\par
\textbf{涉及数据表：}\texttt{visit\_apply}。\par
\textbf{SQL 语句：}
\begin{lstlisting}[style=sqlstyle]
SELECT access_status, COUNT(*) AS status_count
FROM visit_apply
WHERE deleted = 0
GROUP BY access_status
ORDER BY status_count DESC;
\end{lstlisting}
\textbf{查询说明：}该查询与系统页面、后端统计接口或课程设计报告中的业务说明保持一致，能够直接基于当前数据库脚本执行。

\subsection{SQL-21 黑名单风险数量}
\textbf{查询名称：}黑名单风险数量。\par
\textbf{查询目的：}用于黑名单风险分析，按风险级别和状态统计风险数量。\par
\textbf{涉及数据表：}\texttt{blacklist}。\par
\textbf{SQL 语句：}
\begin{lstlisting}[style=sqlstyle]
SELECT level, status, COUNT(*) AS risk_count
FROM blacklist
WHERE deleted = 0
GROUP BY level, status
ORDER BY risk_count DESC;
\end{lstlisting}
\textbf{查询说明：}该查询与系统页面、后端统计接口或课程设计报告中的业务说明保持一致，能够直接基于当前数据库脚本执行。

\subsection{SQL-22 最近操作日志}
\textbf{查询名称：}最近操作日志。\par
\textbf{查询目的：}用于系统日志页，展示最近关键操作记录，支撑系统审计。\par
\textbf{涉及数据表：}\texttt{operation\_log}。\par
\textbf{SQL 语句：}
\begin{lstlisting}[style=sqlstyle]
SELECT operator_name, module_name, operation_type, operation_result, request_url, operation_time
FROM operation_log
WHERE deleted = 0
ORDER BY operation_time DESC
LIMIT 10;
\end{lstlisting}
\textbf{查询说明：}该查询与系统页面、后端统计接口或课程设计报告中的业务说明保持一致，能够直接基于当前数据库脚本执行。

\subsection{建表 SQL 摘要}
数据库建表设计在主键、外键、索引、状态字段和时间字段上保持统一规范。以下摘要展示核心建表语句的结构特征，用于说明物理实现与逻辑结构设计之间的对应关系。
\begin{lstlisting}[style=sqlstyle]
-- 重庆邮电大学智慧访客预约与出入校管理系统
-- MySQL 8 schema script

USE cqupt_visitor_system;
SET NAMES utf8mb4;
SET FOREIGN_KEY_CHECKS = 0;

DROP TABLE IF EXISTS report_record;
DROP TABLE IF EXISTS screenshot_record;
DROP TABLE IF EXISTS dict_item;
DROP TABLE IF EXISTS operation_log;
DROP TABLE IF EXISTS notice;
DROP TABLE IF EXISTS blacklist;
DROP TABLE IF EXISTS access_record;
DROP TABLE IF EXISTS pass_code;
DROP TABLE IF EXISTS approval_record;
DROP TABLE IF EXISTS visitor_companion;
DROP TABLE IF EXISTS visit_apply;
DROP TABLE IF EXISTS visitor_vehicle;
DROP TABLE IF EXISTS sys_role_permission;
DROP TABLE IF EXISTS sys_user_role;
DROP TABLE IF EXISTS sys_user;
DROP TABLE IF EXISTS dict_type;
DROP TABLE IF EXISTS visitor;
DROP TABLE IF EXISTS campus_gate;
DROP TABLE IF EXISTS sys_permission;
DROP TABLE IF EXISTS sys_role;
DROP TABLE IF EXISTS department;

SET FOREIGN_KEY_CHECKS = 1;

CREATE TABLE department (
  id BIGINT NOT NULL AUTO_INCREMENT COMMENT '部门主键',
  parent_id BIGINT NULL COMMENT '上级部门',
  dept_code VARCHAR(50) NOT NULL COMMENT '部门编码',
  dept_name VARCHAR(100) NOT NULL COMMENT '部门名称',
  leader_user_id BIGINT NULL COMMENT '部门负责人用户编号',
  phone VARCHAR(20) NULL COMMENT '联系电话',
  sort_order INT NOT NULL DEFAULT 0 COMMENT '排序号',
  status VARCHAR(30) NOT NULL DEFAULT 'ENABLED' COMMENT '部门状态：ENABLED、DISABLED',
  create_time DATETIME NOT NULL DEFAULT CURRENT_TIMESTAMP COMMENT '创建时间',
  update_time DATETIME NOT NULL DEFAULT CURRENT_TIMESTAMP ON UPDATE CURRENT_TIMESTAMP COMMENT '更新时间',
  deleted TINYINT NOT NULL DEFAULT 0 COMMENT '软删除标记：0未删除，1已删除',
  PRIMARY KEY (id),
  UNIQUE KEY uk_department_code (dept_code),
  KEY idx_department_parent (parent_id),
  KEY idx_department_leader (leader_user_id),
  KEY idx_department_status (status),
  CONSTRAINT fk_department_parent FOREIGN KEY (parent_id) REFERENCES department(id) ON DELETE SET NULL
) ENGINE=InnoDB DEFAULT CHARSET=utf8mb4 COLLATE=utf8mb4_unicode_ci COMMENT='部门表';

CREATE TABLE sys_role (
  id BIGINT NOT NULL AUTO_INCREMENT COMMENT '角色主键',
  role_code VARCHAR(50) NOT NULL COMMENT '角色编码',
  role_name VARCHAR(50) NOT NULL COMMENT '角色名称',
  role_desc VARCHAR(255) NULL COMMENT '角色说明',
  status VARCHAR(30) NOT NULL DEFAULT 'ENABLED' COMMENT '角色状态：ENABLED、DISABLED',
  create_time DATETIME NOT NULL DEFAULT CURRENT_TIMESTAMP COMMENT '创建时间',
  update_time DATETIME NOT NULL DEFAULT CURRENT_TIMESTAMP ON UPDATE CURRENT_TIMESTAMP COMMENT '更新时间',
  deleted TINYINT NOT NULL DEFAULT 0 COMMENT '软删除标记',
  PRIMARY KEY (id),
  UNIQUE KEY uk_sys_role_code (role_code),
  KEY idx_sys_role_status (status)
) ENGINE=InnoDB DEFAULT CHARSET=utf8mb4 COLLATE=utf8mb4_unicode_ci COMMENT='角色表';

CREATE TABLE sys_permission (
  id BIGINT NOT NULL AUTO_INCREMENT COMMENT '权限主键',
  permission_code VARCHAR(100) NOT NULL COMMENT '权限编码',
  permission_name VARCHAR(100) NOT NULL COMMENT '权限名称',
  permission_type VARCHAR(30) NOT NULL DEFAULT 'MENU' COMMENT '权限类型：MENU、BUTTON、API',
  parent_id BIGINT NULL COMMENT '父级权限',
  route_path VARCHAR(200) NULL COMMENT '前端路由',
  component_path VARCHAR(200) NULL COMMENT '前端组件路径',
  api_path VARCHAR(200) NULL COMMENT '后端接口路径',
  sort_order INT NOT NULL DEFAULT 0 COMMENT '排序号',
  status VARCHAR(30) NOT NULL DEFAULT 'ENABLED' COMMENT '权限状态',
  create_time DATETIME NOT NULL DEFAULT CURRENT_TIMESTAMP COMMENT '创建时间',
  update_time DATETIME NOT NULL DEFAULT CURRENT_TIMESTAMP ON UPDATE CURRENT_TIMESTAMP COMMENT '更新时间',
  deleted TINYINT NOT NULL DEFAULT 0 COMMENT '软删除标记',
  PRIMARY KEY (id),
  UNIQUE KEY uk_sys_permission_code (permission_code),
  KEY idx_sys_permission_parent (parent_id),
  KEY idx_sys_permission_type (permission_type),
  CONSTRAINT fk_permission_parent FOREIGN KEY (parent_id) REFERENCES sys_permission(id) ON DELETE SET NULL
) ENGINE=InnoDB DEFAULT CHARSET=utf8mb4 COLLATE=utf8mb4_unicode_ci COMMENT='权限表';

CREATE TABLE campus_gate (
  id BIGINT NOT NULL AUTO_INCREMENT COMMENT '校门主键',
  gate_code VARCHAR(50) NOT NULL COMMENT '校门编码',
  gate_name VARCHAR(100) NOT NULL COMMENT '校门名称',
  gate_location VARCHAR(200) NULL COMMENT '校门位置',
  gate_type VARCHAR(30) NOT NULL DEFAULT 'NORMAL' COMMENT '校门类型：NORMAL、VEHICLE、PEDESTRIAN',
  status VARCHAR(30) NOT NULL DEFAULT 'ENABLED' COMMENT '校门状态',
  remark VARCHAR(255) NULL COMMENT '备注',
  create_time DATETIME NOT NULL DEFAULT CURRENT_TIMESTAMP COMMENT '创建时间',
  update_time DATETIME NOT NULL DEFAULT CURRENT_TIMESTAMP ON UPDATE CURRENT_TIMESTAMP COMMENT '更新时间',
  deleted TINYINT NOT NULL DEFAULT 0 COMMENT '软删除标记',
  PRIMARY KEY (id),
  UNIQUE KEY uk_campus_gate_code (gate_code),
  KEY idx_campus_gate_status (status)
) ENGINE=InnoDB DEFAULT CHARSET=utf8mb4 COLLATE=utf8mb4_unicode_ci COMMENT='校门表';

CREATE TABLE visitor (
  id BIGINT NOT NULL AUTO_INCREMENT COMMENT '访客主键',
  visitor_name VARCHAR(50) NOT NULL COMMENT '访客姓名',
  id_type VARCHAR(30) NOT NULL DEFAULT 'ID_CARD' COMMENT '证件类型',
  id_number VARCHAR(50) NOT NULL COMMENT '证件号码',
  phone VARCHAR(20) NOT NULL COMMENT '手机号',
  company VARCHAR(100) NULL COMMENT '单位',
  gender VARCHAR(10) NULL COMMENT '性别',
  visitor_level VARCHAR(30) NOT NULL DEFAULT 'NORMAL' COMMENT '访客等级：NORMAL、VIP、RISK',
  status VARCHAR(30) NOT NULL DEFAULT 'NORMAL' COMMENT '访客状态',
  create_time DATETIME NOT NULL DEFAULT CURRENT_TIMESTAMP COMMENT '创建时间',
  update_time DATETIME NOT NULL DEFAULT CURRENT_TIMESTAMP ON UPDATE CURRENT_TIMESTAMP COMMENT '更新时间',
  deleted TINYINT NOT NULL DEFAULT 0 COMMENT '软删除标记',
  PRIMARY KEY (id),
  UNIQUE KEY uk_visitor_id_number (id_type, id_number),
  KEY idx_visitor_phone (phone),
  KEY idx_visitor_status (status)
) ENGINE=InnoDB DEFAULT CHARSET=utf8mb4 COLLATE=utf8mb4_unicode_ci COMMENT='访客表';

CREATE TABLE dict_type (
  id BIGINT NOT NULL AUTO_INCREMENT COMMENT '字典类型主键',
  type_code VARCHAR(50) NOT NULL COMMENT '字典类型编码',
  type_name VARCHAR(100) NOT NULL COMMENT '字典类型名称',
  status VARCHAR(30) NOT NULL DEFAULT 'ENABLED' COMMENT '状态',
  remark VARCHAR(255) NULL COMMENT '备注',
  create_time DATETIME NOT NULL DEFAULT CURRENT_TIMESTAMP COMMENT '创建时间',
  update_time DATETIME NOT NULL DEFAULT CURRENT_TIMESTAMP ON UPDATE CURRENT_TIMESTAMP COMMENT '更新时间',
  deleted TINYINT NOT NULL DEFAULT 0 COMMENT '软删除标记',
  PRIMARY KEY (id),
  UNIQUE KEY uk_dict_type_code (type_code),
  KEY idx_dict_type_status (status)
) ENGINE=InnoDB DEFAULT CHARSET=utf8mb4 COLLATE=utf8mb4_unicode_ci COMMENT='字典类型表';

CREATE TABLE sys_user (
  id BIGINT NOT NULL AUTO_INCREMENT COMMENT '用户主键',
  username VARCHAR(50) NOT NULL COMMENT '登录账号',
  password_hash VARCHAR(255) NOT NULL COMMENT '加密后的密码',
  real_name VARCHAR(50) NOT NULL COMMENT '真实姓名',
  phone VARCHAR(20) NOT NULL COMMENT '手机号',
  email VARCHAR(100) NULL COMMENT '邮箱',
  department_id BIGINT NULL COMMENT '所属部门',
  user_type VARCHAR(30) NOT NULL COMMENT '用户类型：HOST、APPROVER、GUARD、ADMIN、MANAGER',
  status VARCHAR(30) NOT NULL DEFAULT 'ENABLED' COMMENT '账号状态',
  last_login_time DATETIME NULL COMMENT '最近登录时间',
  create_time DATETIME NOT NULL DEFAULT CURRENT_TIMESTAMP COMMENT '创建时间',
  update_time DATETIME NOT NULL DEFAULT CURRENT_TIMESTAMP ON UPDATE CURRENT_TIMESTAMP COMMENT '更新时间',
  deleted TINYINT NOT NULL DEFAULT 0 COMMENT '软删除标记',
  PRIMARY KEY (id),
  UNIQUE KEY uk_sys_user_username (username),
  KEY idx_sys_user_phone (phone),
  KEY idx_sys_user_department (department_id),
  KEY idx_sys_user_status (status),
  CONSTRAINT fk_user_department FOREIGN KEY (department_id) REFERENCES department(id) ON DELETE SET NULL
) ENGINE=InnoDB DEFAULT CHARSET=utf8mb4 COLLATE=utf8mb4_unicode_ci COMMENT='系统用户表';

ALTER TABLE department
  ADD CONSTRAINT fk_department_leader FOREIGN KEY (leader_user_id) REFERENCES sys_user(id) ON DELETE SET NULL;

CREATE TABLE sys_user_role (
  id BIGINT NOT NULL AUTO_INCREMENT COMMENT '关联主键',
  user_id BIGINT NOT NULL COMMENT '用户编号',
  role_id BIGINT NOT NULL COMMENT '角色编号',
  create_time DATETIME NOT NULL DEFAULT CURRENT_TIMESTAMP COMMENT '创建时间',
  update_time DATETIME NOT NULL DEFAULT CURRENT_TIMESTAMP ON UPDATE CURRENT_TIMESTAMP COMMENT '更新时间',
  deleted TINYINT NOT NULL DEFAULT 0 COMMENT '软删除标记',
  PRIMARY KEY (id),
  UNIQUE KEY uk_user_role (user_id, role_id),
  KEY idx_user_role_role (role_id),
  CONSTRAINT fk_user_role_user FOREIGN KEY (user_id) REFERENCES sys_user(id) ON DELETE CASCADE,
  CONSTRAINT fk_user_role_role FOREIGN KEY (role_id) REFERENCES sys_role(id) ON DELETE CASCADE
) ENGINE=InnoDB DEFAULT CHARSET=utf8mb4 COLLATE=utf8mb4_unicode_ci COMMENT='用户角色关联表';

CREATE TABLE sys_role_permission (
  id BIGINT NOT NULL AUTO_INCREMENT COMMENT '关联主键',
  role_id BIGINT NOT NULL COMMENT '角色编号',
  permission_id BIGINT NOT NULL COMMENT '权限编号',
  create_time DATETIME NOT NULL DEFAULT CURRENT_TIMESTAMP COMMENT '创建时间',
  update_time DATETIME NOT NULL DEFAULT CURRENT_TIMESTAMP ON UPDATE CURRENT_TIMESTAMP COMMENT '更新时间',
\end{lstlisting}
